\documentclass[a4paper,11pt]{article}

\usepackage[T1]{fontenc}
\usepackage[utf8]{inputenc}
\usepackage{lmodern}
\usepackage{microtype}
\usepackage[a4paper,margin=1in,headheight=15pt,headsep=14pt]{geometry}
\usepackage{amsmath,amssymb}
\usepackage{graphicx}
\usepackage{booktabs}
\usepackage[table,svgnames]{xcolor}
\usepackage[most]{tcolorbox}
\tcbuselibrary{skins,breakable}

\usepackage{pifont}
\IfFileExists{bbding.sty}{%
  \usepackage{bbding}%
  \newcommand{\glyphHand}{\Pointinghand}%
  \newcommand{\glyphPencil}{\Pencil}%
}{%
  \newcommand{\glyphHand}{\ding{43}}%
  \newcommand{\glyphPencil}{\ding{46}}%
}

\usepackage{enumitem}
\setlist{leftmargin=1.5em,topsep=4pt,itemsep=3pt}

\newlist{steps}{enumerate}{1}
\setlist[steps]{label=\textbf{Step \arabic*.},
  leftmargin=4.4em, labelwidth=3.8em, labelsep=0.6em, labelindent=0pt,
  align=left, topsep=5pt, itemsep=6pt}
\usepackage{parskip}

\usepackage{fancyhdr}
\usepackage{titlesec}
\usepackage[hidelinks]{hyperref}
\usepackage{xurl}
\hypersetup{breaklinks=true}

\definecolor{navy}{HTML}{21355E}
\definecolor{navyrule}{HTML}{21355E}
\definecolor{inktext}{HTML}{1A202C}
\definecolor{mutedtext}{HTML}{6B7280}

\definecolor{intuFrame}{HTML}{2C5AA0}   \definecolor{intuFill}{HTML}{F4F7FC}
\definecolor{everyFrame}{HTML}{8A5A1E}  \definecolor{everyFill}{HTML}{FBF1E3}
\definecolor{workFrame}{HTML}{2E7D52}   \definecolor{workFill}{HTML}{F1F8F3}
\definecolor{watchFrame}{HTML}{B23A48}  \definecolor{watchFill}{HTML}{FBF1F2}
\definecolor{keyFrame}{HTML}{6A4C93}    \definecolor{keyFill}{HTML}{F4F0F9}
\definecolor{waitFrame}{HTML}{0F766E}   \definecolor{waitFill}{HTML}{EEF7F5}

\color{inktext}
\hypersetup{linkcolor=navy,urlcolor=navy,citecolor=navy}

\newcommand{\CourseShort}{DNN}
\newcommand{\SessionNo}{15}
\newcommand{\LectureTitle}{Regularization of Deep Models}

\pagestyle{fancy}
\fancyhf{}
\fancyhead[L]{\footnotesize\color{mutedtext}\textsf{\CourseShort}}
\fancyhead[R]{\footnotesize\color{mutedtext}\textsf{Session \SessionNo\ \textperiodcentered\ \LectureTitle}}
\fancyfoot[L]{\footnotesize\color{mutedtext}\textsf{WILP, BITS Pilani}}
\fancyfoot[C]{\footnotesize\color{mutedtext}\thepage}
\renewcommand{\headrulewidth}{0.4pt}
\renewcommand{\footrulewidth}{0pt}
\renewcommand{\headrule}{\hbox to\headwidth{\color{navyrule!55}\leaders\hrule height \headrulewidth\hfill}}

\titleformat{\section}
  {\sffamily\Large\bfseries\color{navy}}
  {\thesection}
  {0.8em}
  {}
  [{\vspace{2pt}\color{navy!70}\titlerule[0.8pt]}]
\titleformat{\subsection}
  {\sffamily\large\bfseries\color{navy!92}}
  {\thesubsection}{0.7em}{}
\titlespacing*{\section}{0pt}{16pt}{8pt}
\titlespacing*{\subsection}{0pt}{12pt}{4pt}

\newif\ifmlkfirstsec \mlkfirstsectrue
\newcommand{\sectionbreak}{\ifmlkfirstsec\mlkfirstsecfalse\else\clearpage\fi}

\newcommand{\secsub}[1]{%
  \noindent{\sffamily\bfseries\itshape\color{navy!85}\large #1}\par\medskip}

\newcommand{\flipanswer}[1]{\rotatebox[origin=c]{180}{\parbox{0.92\linewidth}{\small #1}}}

\tcbset{
  housecallout/.style 2 args={
    enhanced, breakable,
    colback=#2, colframe=#1,
    boxrule=0.6pt, arc=2.4pt, outer arc=2.4pt,
    left=10pt, right=10pt, top=9pt, bottom=9pt,
    before skip=11pt, after skip=11pt,
    fontupper=\normalsize,
    coltitle=white,
    fonttitle=\bfseries\sffamily\small,
    attach boxed title to top left={xshift=6mm,yshift=-3mm},
    boxed title style={
      colback=#1, colframe=#1,
      rounded corners, arc=3pt,
      boxrule=0pt, left=6pt, right=6pt, top=2pt, bottom=2pt
    }
  }
}

\newtcolorbox{intuition}{housecallout={intuFrame}{intuFill}, title={\raisebox{-0.05ex}{\glyphHand}\,\ The intuition}}
\newtcolorbox{everyday}{housecallout={everyFrame}{everyFill}, title={\raisebox{-0.05ex}{$\bigstar$}\,\ Everyday picture}}
\newtcolorbox{worked}[1]{housecallout={workFrame}{workFill}, title={\raisebox{-0.05ex}{\glyphPencil}\,\ Worked example: #1}}
\newtcolorbox{watchout}{housecallout={watchFrame}{watchFill}, title={\raisebox{-0.05ex}{\ding{55}}\,\ Watch out}}
\newtcolorbox{keytake}{housecallout={keyFrame}{keyFill}, title={\raisebox{-0.05ex}{\ding{51}}\,\ Key takeaway}}
\newtcolorbox{waitwhat}{housecallout={waitFrame}{waitFill}, title={\raisebox{-0.05ex}{\ding{93}}\,\ Wait --- really?}}

\newcommand{\bitswordmark}{{\sffamily\bfseries\color{white}\fontsize{12.5}{15}\selectfont WILP, BITS Pilani}}
\newcommand{\bitslogochip}{}

\newcommand{\titlebanner}[4]{%
  \begin{tcolorbox}[
    enhanced, breakable=false,
    colback=navy, colframe=navy,
    boxrule=0pt, arc=3pt, outer arc=3pt,
    left=16pt, right=16pt, top=16pt, bottom=16pt,
    before skip=2pt, after skip=14pt,
    width=\linewidth ]
    \noindent
    \begin{minipage}[c]{0.72\linewidth}\bitswordmark\end{minipage}%
    \hfill
    \begin{minipage}[c]{0.26\linewidth}\raggedleft\bitslogochip\end{minipage}\par
    \vspace{9pt}
    {\sffamily\footnotesize\bfseries\color{white!85}%
      \MakeUppercase{#1}}\par
    \vspace{6pt}
    {\sffamily\bfseries\fontsize{26}{30}\selectfont\color{white} #2\par}
    \vspace{5pt}
    {\sffamily\large\color{white!88} #3\par}
    \vspace{9pt}
    {\color{white!55}\rule{\linewidth}{0.4pt}}\par
    \vspace{6pt}
    {\sffamily\footnotesize\color{white!82} #4\par}
  \end{tcolorbox}%
}

\newcommand{\housefig}[2]{%
  \begin{figure}[htbp]
    \centering
    \includegraphics[width=\linewidth]{#1}%
    \caption{#2}
  \end{figure}%
}

\newcommand{\vocab}[1]{\textbf{\color{navy}#1}}

\begin{document}

\thispagestyle{fancy}

\titlebanner
  {Deep Neural Networks \textperiodcentered\ Module 11 Companion}
  {Regularization of Deep Models}
  {Controlling Overfitting through Norm Penalties, Dropout, Augmentation, and Early Stopping}
  {Session 15 (23rd August 2026) \textperiodcentered\ Read alongside the lecture slides \textperiodcentered\ Every concept explained in full}

\smallskip
\noindent{\small\color{mutedtext}%
Prepared for the Work Integrated Learning Programmes (WILP), BITS Pilani.}
\smallskip

\section{The Generalization Problem in Deep Learning}
\secsub{Overfitting occurs when a high-capacity neural network fits noise in the training set rather than the underlying data distribution.}

Deep neural networks possess millions of parameters, giving them immense capacity.
Without constraints, the model memorizes training samples, achieving near-zero training loss while performing poorly on unseen validation data.
\vocab{Regularization} is any modification made to a learning algorithm intended to reduce its generalization error but not its training error.

\section{Parameter Norm Penalties (L1 and L2 Regularization)}
\secsub{Adding a penalty term to the objective function restricts model capacity by shrinking parameter magnitudes.}

The regularized objective function $\tilde{J}(w; X, y)$ adds a penalty $\Omega(w)$ scaled by hyperparameter $\alpha \ge 0$:
\begin{equation}
  \tilde{J}(w; X, y) = J(w; X, y) + \alpha \Omega(w)
\end{equation}

\subsection{L2 Regularization (Weight Decay)}
In \vocab{L2 regularization}, $\Omega(w) = \frac{1}{2} \|w\|_2^2 = \frac{1}{2} \sum_i w_i^2$.
The gradient of the regularized loss is:
\begin{equation}
  \nabla_w \tilde{J}(w) = \nabla_w J(w) + \alpha w
\end{equation}
The gradient descent update step becomes:
\begin{equation}
  w_{t+1} = w_t - \eta (\nabla_w J(w_t) + \alpha w_t) = (1 - \eta \alpha) w_t - \eta \nabla_w J(w_t)
\end{equation}
Notice that before taking the gradient step, $w_t$ is multiplied by $(1 - \eta \alpha) < 1$. This is why L2 regularization is known as \vocab{weight decay}.

\subsection{L1 Regularization (Lasso & Sparsity)}
In \vocab{L1 regularization}, $\Omega(w) = \|w\|_1 = \sum_i |w_i|$.
The gradient is:
\begin{equation}
  \nabla_w \tilde{J}(w) = \nabla_w J(w) + \alpha \text{sign}(w)
\end{equation}
L1 regularization penalizes weights by a constant rate $\alpha$ regardless of their magnitude, forcing uninformative feature weights strictly to $0.0$, producing \vocab{sparse weight matrices}.

\housefig{figures/regularization_geometry.png}{Geometric comparison of L1 (Lasso) and L2 (Weight Decay) constraint regions.}

\section{Dropout, Data Augmentation, and Early Stopping}
\secsub{Implicit regularization techniques prevent co-adaptation and enlarge effective dataset sizes.}

\subsection{Dropout}
During training, \vocab{Dropout} randomly deactivates each hidden unit with probability $p$ (e.g. $p=0.5$).
This forces units to learn independent, robust representations rather than co-adapting with specific neighbors.
During evaluation, dropout is disabled, and activations are scaled by $(1-p)$ (or inverted dropout scales by $1/(1-p)$ during training).

\subsection{Data Augmentation & Early Stopping}
Data augmentation creates synthetic training samples using domain transformations (e.g., random flips, rotations, color jittering in vision tasks).
Early stopping monitors validation set performance across training epochs and stops training when validation loss stops improving for $P$ consecutive epochs (\vocab{patience}).

\section{Weight Initialization & Label Smoothing}
\secsub{Proper initialization prevents vanishing or exploding gradients at the start of training.}

Xavier (Glorot) initialization draws weights from $\mathcal{N}(0, \sigma^2)$ with $\sigma = \sqrt{\frac{2}{n_{in} + n_{out}}}$, ideal for Sigmoid/Tanh activations.
Kaiming (He) initialization sets $\sigma = \sqrt{\frac{2}{n_{in}}}$, specially designed for ReLU networks.
Label smoothing replaces hard binary targets $y \in \{0, 1\}$ with smoothed targets $y_{smooth} = (1 - \epsilon) y + \frac{\epsilon}{K}$, preventing the model from becoming overly confident in its predictions.

\section*{Self-Test \& Summary}
\addcontentsline{toc}{section}{Self-Test \& Summary}

\subsection*{Self-Test Questions}
\begin{enumerate}
  \item Why does L1 regularization create sparse weights while L2 does not?
        \flipanswer{L1's diamond constraint region has sharp corners along coordinate axes where loss contours hit first, setting weights to exactly 0.}
  \item Why is inverted dropout preferred in modern frameworks?
        \flipanswer{It scales activations by $1/(1-p)$ during training so evaluation code runs unchanged without scaling.}
\end{enumerate}

\end{document}
